\section{Conclusion and Recommendations}
\label{sec:conclusion}

\subsection{Summary of Key Findings}

This research addresses the critical challenge of optimizing humanitarian aid delivery routes under road network uncertainty, focusing on vaccine and life-saving supply distribution in Somalia. Our work makes novel contributions to stochastic vehicle routing and humanitarian logistics literature.

\subsubsection{Research Contributions}

\textbf{1. Stochastic VRP with Reliability Constraints.}
We developed a chance-constrained stochastic programming framework that explicitly models road travel time uncertainty using probability distributions fitted to 5 years of Somalia road network data. Unlike traditional SVRP models focusing on expected cost minimization, our approach treats reliability as a hard constraint with threshold $\alpha = 0.85$, ensuring minimum service standards critical for perishable humanitarian supplies.

\textbf{2. Progressive Modeling Framework.}
Our Model 0 $\to$ Model 1 $\to$ Model 2 $\to$ Model 3 progression demonstrates systematic problem understanding:
\begin{itemize}
\item Model 0 (deterministic VRP) establishes performance baseline
\item Model 1 (chance-constrained SVRP) demonstrates stochastic optimization benefits
\item Model 2 (two-stage stochastic program) improves computational efficiency via Benders decomposition
\item Model 3 (multi-depot SVRP) extends to coordinated multi-warehouse strategies
\item Model 4 (multi-objective) identifies efficiency-reliability-cost trade-offs
\end{itemize}

\textbf{3. Hybrid Solution Algorithm.}
We designed a hybrid solution approach combining exact methods (CPLEX) for validation and Adaptive Large Neighborhood Search (ALNS) metaheuristic for realistic instances. Our implementation achieves near-optimal solutions (gap $< 4\%$) within 3-4 minutes for instances up to 60 customers, enabling practical field implementation.

\textbf{4. Multi-Warehouse Coordination.}
We developed a novel clustering-first, routing-second approach for multi-depot stochastic VRP, demonstrating 31\% improvement in delivery efficiency compared to centralized single-warehouse distribution.

\subsubsection{Empirical Results}

Our computational experiments on Somalia road network data demonstrate:

\begin{table}[H]
\centering
\caption{Summary of Key Results}
\begin{tabular}{lcc}
\toprule
Metric & Deterministic Baseline & Stochastic Optimization \\
\midrule
Expected Delivery Time & 47.3 hours & 36.5 hours (-23\%) \\
On-Time Delivery Rate & 67\% & 89\% (+22 percentage points) \\
Operational Cost & \$2,340 & \$2,120 (-9\%) \\
Multi-Warehouse Improvement & N/A & 31\% reduction \\
\bottomrule
\end{tabular}
\end{table}

These improvements are substantial: 23\% reduction in delivery time translates to faster life-saving supply delivery; 22 percentage point reliability improvement means 89 out of 100 deliveries arrive on time versus 67 under deterministic routing.

\subsection{Practical Recommendations for Humanitarian Organizations}

\subsubsection{Recommendation 1: Adopt Stochastic Routing with Reliability Thresholds}

\textbf{Action:} Replace deterministic routing with stochastic optimization using reliability thresholds $\alpha \in [0.85, 0.90]$.

\textbf{Rationale:}
\begin{itemize}
\item 23\% reduction in delivery time
\item 22 percentage point improvement in on-time delivery rate
\item Minimal computational overhead (3-4 minutes runtime)
\end{itemize}

\textbf{Implementation Steps:}
\begin{enumerate}
\item Collect 2-3 years of historical road travel time data
\item Fit probability distributions (lognormal, Weibull) to each road segment
\item Implement ALNS algorithm (open-source code provided in Appendix C)
\item Generate daily route plans with reliability $\alpha = 0.85$
\item Validate with 1000-scenario out-of-sample testing
\end{enumerate}

\subsubsection{Recommendation 2: Implement Multi-Warehouse Coordination}

\textbf{Action:} Transition from centralized single-warehouse distribution to coordinated multi-depot network.

\textbf{Rationale:}
\begin{itemize}
\item 31\% improvement in delivery efficiency
\item Better geographic coverage (closer warehouses to beneficiaries)
\item Increased resilience (single warehouse failure less catastrophic)
\end{itemize}

\textbf{Implementation Steps:}
\begin{enumerate}
\item Assess current warehouse network and identify optimal locations for new warehouses
\item Implement reliability-weighted customer clustering (Algorithm 3)
\item Establish inter-warehouse coordination protocols
\item Train warehouse managers on coordinated operations
\item Monitor performance and iteratively optimize
\end{enumerate}

\subsubsection{Recommendation 3: Establish Data Collection Infrastructure}

\textbf{Action:} Install GPS tracking systems on delivery vehicles for continuous data collection.

\textbf{Rationale:}
\begin{itemize}
\item Better data $\rightarrow$ more accurate probability distributions $\rightarrow$ better routes
\item Real-time monitoring enables dynamic re-routing for extreme events
\item Continuous model improvement via online learning
\end{itemize}

\textbf{Implementation Steps:}
\begin{enumerate}
\item Purchase GPS trackers (\$50-\$100 per vehicle)
\item Develop data pipeline (GPS $\rightarrow$ database $\rightarrow$ probability model)
\item Retrain probability models monthly with new data
\item Monitor distribution shift (concept drift detection)
\end{enumerate}

\subsubsection{Recommendation 4: Develop Decision Support Tool}

\textbf{Action:} Create user-friendly software interface for non-technical users.

\textbf{Rationale:}
\begin{itemize}
\item Lowers implementation barrier (no operations research expertise required)
\item Standardizes routing decisions across organization
\item Enables scenario analysis (what-if planning for emergencies)
\end{itemize}

\textbf{Features:}
\begin{itemize}
\item GIS map visualization of routes
\item One-click route generation
\item Reliability dashboard (predicted vs. actual on-time rates)
\item Sensitivity analysis tools
\end{itemize}

\subsubsection{Recommendation 5: Maintain Contingency Plans for Extreme Events}

\textbf{Action:} Develop manual override protocols for extreme events (flash floods, major security incidents) where model assumptions fail.

\textbf{Rationale:}
\begin{itemize}
\item Model assumes stationary distributions (violated during extreme events)
\item Communication infrastructure may be disrupted (cannot send new routes)
\item Human judgment needed for unprecedented situations
\end{itemize}

\textbf{Protocol:}
\begin{enumerate}
\item Define ``extreme event'' triggers (e.g., travel time increase $> 50\%$)
\item Pre-plan backup routes for critical corridors
\item Stockpile emergency supplies at strategic locations
\item Establish satellite communication backup
\end{enumerate}

\subsection{Recommendations for Humanitarian Logistics Policy}

\subsubsection{Policy Recommendation 1: Data Sharing Consortium}

\textbf{Proposal:} Establish regional humanitarian logistics data sharing consortium.

\textbf{Benefits:}
\begin{itemize}
\item Pool road network data across organizations (economies of scale)
\item Improve probability distribution estimates (more data)
\item Standardize data formats and quality metrics
\end{itemize}

\textbf{Implementation:} UN OCHA (Office for the Coordination of Humanitarian Affairs) could host centralized database with anonymized data sharing agreements.

\subsubsection{Policy Recommendation 2: Pre-Positioned Strategic Stocks}

\textbf{Proposal:} Pre-position emergency supplies at geographically distributed warehouses.

\textbf{Rationale:} Our multi-warehouse results show 31\% efficiency improvement. Pre-positioning reduces last-mile delivery distance.

\textbf{Implementation:} World Food Programme and UNICEF could jointly fund warehouse construction and inventory management.

\subsubsection{Policy Recommendation 3: Investment in Road Infrastructure}

\textbf{Proposal:} Prioritize road infrastructure investments in high-impact corridors.

\textbf{Rationale:} Our sensitivity analysis identifies high-CV road segments as bottlenecks. Paving these segments yields disproportionate reliability improvements.

\textbf{Investment Strategy:}
\begin{itemize}
\item Focus on Southern-Central Corridor (connects Kismayo to Mogadishu)
\item Paving high-CV dirt tracks reduces CV from 0.46 to 0.21
\item Cost-benefit analysis shows 400\% ROI (from reduced logistics costs)
\end{itemize}

\subsection{Limitations and Future Research Directions}

\subsubsection{Limitations}

As discussed in Section \ref{sec:evaluation}, our model has limitations:
\begin{itemize}
\item Static planning (no dynamic re-optimization)
\item Independence assumption (travel times may be correlated)
\item Single-period model (no multi-day inventory optimization)
\item Somalia-specific (generalizability to other contexts untested)
\end{itemize}

\subsubsection{Future Research Direction 1: Dynamic Stochastic VRP}

\textbf{Research Question:} How can we incorporate real-time information (weather forecasts, security alerts) for dynamic re-optimization?

\textbf{Approach:}
\begin{itemize}
\item Rolling horizon optimization with model predictive control
\item Reserve computational budget for rapid re-optimization
\item Machine learning for disruption prediction
\end{itemize}

\subsubsection{Future Research Direction 2: Multi-Period Inventory-Routing}

\textbf{Research Question:} How can we jointly optimize inventory decisions and routing over multi-day planning horizons?

\textbf{Approach:}
\begin{itemize}
\item Extend to Inventory-Routing Problem (IRP) with uncertainty
\item Coordinate inventory replenishment with delivery routes
\item Balance holding costs vs. delivery costs
\end{itemize}

\subsubsection{Future Research Direction 3: Correlated Travel Times}

\textbf{Research Question:} How can we model correlated travel times while maintaining computational tractability?

\textbf{Approach:}
\begin{itemize}
\item Copula models for dependency structure
\item Scenario reduction techniques to keep problem size manageable
\item Analytical bounds on correlation impact
\end{itemize}

\subsubsection{Future Research Direction 4: Humanitarian Behavior Integration}

\textbf{Research Question:} How do human factors (driver behavior, beneficiary availability) affect routing optimization?

\textbf{Approach:}
\begin{itemize}
\item Behavioral experiments to understand decision-making under uncertainty
\item Incorporate behavioral constraints into optimization models
\item Design nudges to improve plan adherence
\end{itemize}

\subsubsection{Future Research Direction 5: Generalization to Other Contexts}

\textbf{Research Question:} How well does our approach generalize to other humanitarian contexts (natural disasters, refugee crises)?

\textbf{Approach:}
\begin{itemize}
\item Test model on datasets from different regions (e.g., Syria, Yemen)
\item Identify context-specific modifications needed
\item Develop general framework with customizable parameters
\end{itemize}

\subsection{Concluding Remarks}

Humanitarian logistics plays a critical role in delivering life-saving supplies to populations in need. This research demonstrates that rigorous stochastic optimization can significantly improve delivery efficiency and reliability in uncertain environments.

Our key message to humanitarian practitioners: \textbf{Embrace uncertainty, don't ignore it.} Traditional deterministic routing fails when road conditions are unpredictable. Stochastic optimization with reliability constraints provides robust solutions that perform well under uncertainty, not just on average.

To the operations research community, we offer a rigorous stochastic VRP formulation with reliability constraints, effective hybrid solution algorithms, and comprehensive validation using real-world data. The 23\% reduction in delivery time and 22 percentage point reliability improvement demonstrate the practical value of our approach.

Most importantly, this research has real-world impact: faster delivery means vaccines reach children before they expire; reliable delivery means health centers can plan operations with confidence; efficient delivery means limited resources serve more people. In humanitarian logistics, better algorithms literally save lives.

\subsection{Final Words}

We developed this model with the Somali context in mind, but the principles apply broadly: understand uncertainty, model it rigorously, optimize with practical constraints in mind, and validate with real data. Humanitarian logistics deserves our best analytical work—because the stakes are life and death.

\textit{Control Number: IMMC26601951}

\textit{Date: March 2026}
